\chapter{Estado del arte}

\section{Introducción y alcance del capítulo}

Este capítulo mapea el trasfondo técnico de un defensor de ciberseguridad basado en aprendizaje por refuerzo que opera sobre observaciones de red a nivel de flujo. El alcance es deliberadamente más amplio que una comparación de clasificadores binarios. La tesis combina detección de intrusiones en red, representación de características de flujo, conjuntos de datos de referencia públicos, líneas base de aprendizaje supervisado, una formulación de conjunto de datos como entorno al estilo Gymnasium, QRDQN, preprocesamiento con control de fugas, comprobaciones de validación más estrictas, experimentos de eficiencia de datos, y validación planificada o preferida sobre tráfico capturado por el operador en un laboratorio doméstico cerrado y no adversarial.

El capítulo sigue por tanto un embudo. Primero revisa la detección de intrusiones en red y la representación del tráfico, luego sitúa a CICIDS2017 entre los conjuntos de datos públicos, después discute las líneas base de aprendizaje supervisado y aprendizaje profundo, y finalmente enmarca la formulación de RL tanto dentro de la teoría del aprendizaje por refuerzo como en trabajos previos de RL/DRL para detección de intrusiones y defensa cibernética. Las últimas secciones se centran en los riesgos metodológicos, los protocolos de evaluación, la evidencia sobre escala de entrenamiento y la brecha de investigación que aborda esta tesis.

La afirmación central es cautelosa. El aprendizaje por refuerzo para la detección de intrusiones ya existe, y los benchmarks públicos de NIDS son ampliamente utilizados \parencite{Yang2024DRLNIDSSurvey,Kheddar2025RLIDSReview,Gueriani2024DRLIoTIDSSurvey}. Esta tesis no pretende introducir el primer IDS basado en RL, demostrar que QRDQN es el estado del arte para NIDS, ni demostrar una prevención en línea lista para producción. Estudia un pipeline de decisión reproducible, offline, a nivel de flujo, con acciones \texttt{PERMIT}/\texttt{BLOCK}, y evalúa cómo se comporta dicha formulación bajo benchmarks internos controlados, comparación con líneas base y validación progresivamente más estricta.

La revisión también separa tres niveles que con frecuencia se confunden. El primero es la monitorización: recopilar paquetes, flujos, registros o eventos derivados y decidir si son sospechosos. El segundo es la predicción: entrenar un modelo que asigne representaciones de tráfico a etiquetas benignas o maliciosas. El tercero es el control: tomar acciones que modifiquen el estado futuro de la red. Esta tesis trabaja principalmente en los dos primeros niveles. Utiliza el lenguaje de \texttt{PERMIT} y \texttt{BLOCK} porque el entorno expone una decisión de seguridad binaria, pero la implementación actual no es un plano de control en línea. Esa distinción condiciona todas las afirmaciones posteriores.

\section{Detección de intrusiones y monitorización de seguridad en red}

Los sistemas de detección de intrusiones monitorizan la actividad de hosts o redes en busca de evidencias de ataques, usos indebidos o violaciones de políticas \parencite{ScarfoneMell2007}. Un sistema de detección de intrusiones en host observa endpoints, eventos del sistema operativo, registros, procesos, cambios en la integridad de archivos o llamadas al sistema. Un sistema de detección de intrusiones en red observa el tráfico que atraviesa enlaces, taps, puertos de mirror de switches, brokers de paquetes, interfaces de red virtual o colectores de flujos. Ambas perspectivas son complementarias: el HIDS puede ver el contexto local de procesos y archivos que los sensores de red no pueden, mientras que el NIDS puede inspeccionar movimientos laterales, escaneos, tráfico de mando y control, y patrones de comunicación agregados sin instalar un agente en cada host.

También es necesario distinguir el IDS de los sistemas de prevención de intrusiones. Un IDS detecta y notifica principalmente actividad sospechosa. Su salida puede ser una alerta, un evento de registro, un registro de flujo enriquecido, un ticket, un evento SIEM o un artefacto forense. Un IPS se posiciona en línea y puede bloquear, descartar, modificar, limitar la tasa, reiniciar o redirigir el tráfico \parencite{ScarfoneMell2007}. En la práctica, los productos y arquitecturas frecuentemente difuminan esta frontera porque un motor de detección puede alimentar reglas de firewall, acciones EDR, playbooks SOAR o actualizaciones de política SDN. Para la evaluación, sin embargo, la distinción importa. Un detector puede tolerar un perfil de falsos positivos diferente al de un bloqueador en línea, dado que las alertas falsas principalmente afectan a la carga de trabajo del analista, mientras que los bloqueos falsos pueden interrumpir tráfico legítimo de negocio.

Las taxonomías clásicas de IDS distinguen enfoques basados en firmas, basados en anomalías, basados en especificaciones e híbridos \parencite{ScarfoneMell2007}. La detección basada en firmas coteja patrones maliciosos conocidos: secuencias de bytes, predicados de reglas, condiciones de protocolo, cadenas de comandos, rastros de exploits o combinaciones de comportamiento conocido. Su fortaleza es la precisión y la explicabilidad ante amenazas conocidas. Su debilidad es la dependencia del mantenimiento de reglas y la escasa cobertura de ataques nuevos, comportamiento polimórfico, cargas útiles cifradas y variaciones específicas del entorno.

La detección basada en anomalías modela el comportamiento esperado y señala las desviaciones \parencite{SommerPaxson2010ClosedWorld,Ring2019DatasetSurvey}. Esta es la familia más naturalmente asociada con el NIDS basado en ML, pero la palabra anomalía no debe interpretarse como sinónimo automático de aprendizaje no supervisado. Muchos estudios públicos de NIDS son clasificadores supervisados entrenados sobre flujos etiquetados como benignos y de ataque, incluso cuando la motivación más amplia es la detección de anomalías. El reto operativo es que el tráfico benigno es amplio y no estacionario, la prevalencia de ataques es baja, las etiquetas son incompletas, y los cambios en la actividad de negocio pueden parecer anómalos sin ser maliciosos.

La detección basada en especificaciones codifica el comportamiento permitido o esperado, por ejemplo máquinas de estado de protocolo, secuencias de comandos permitidas o flujos de trabajo de control industrial restringidos. Puede ser potente cuando el sistema monitorizado tiene un comportamiento estrecho y estable, pero requiere conocimiento del dominio y mantenimiento. Los sistemas híbridos combinan estas ideas, como reglas de firmas para exploits conocidos, puntuaciones estadísticas de anomalías para comportamiento de flujo inusual, y telemetría de host para confirmación en el endpoint. Los programas modernos de monitorización raramente dependen de un único paradigma de detección.

La generación de alertas y la respuesta activa forman un eje separado. Un sistema puede limitarse a producir alertas para su posterior clasificación, puede enriquecer las alertas con evidencia contextual, puede recomendar pasos de respuesta o puede desencadenar contención activa. Esta tesis se sitúa deliberadamente en el extremo conservador de ese eje. Estudia un modelo de decisión offline similar a un NIDS sobre filas de flujo extraídas. El modelo emite predicciones \texttt{PERMIT}/\texttt{BLOCK} durante los experimentos, pero no descarta paquetes, no reescribe políticas de firewall, no actualiza reglas SDN, no pone en cuarentena hosts ni opera en un bucle de retroalimentación en vivo. La bibliografía de NIDS es, por tanto, relevante, pero las afirmaciones sobre IPS en producción quedan fuera del alcance soportado.

\section{Representación del tráfico de red}

El tráfico de red puede representarse como paquetes, bytes de carga útil, sesiones, conexiones o flujos. Las representaciones a nivel de paquete preservan la temporización, las cabeceras, los tamaños, las direcciones, los flags y, en ocasiones, parte de la carga útil por paquete. Son adecuadas para el análisis de protocolos a grano fino y el modelado de secuencias, pero resultan costosas de almacenar y procesar en redes de alta velocidad. Las representaciones a nivel de carga útil retienen el contenido de la aplicación y pueden identificar exploits o firmas de malware invisibles en los metadatos. Su utilidad está limitada por el cifrado, las restricciones de privacidad, las restricciones legales y el coste de CPU.

Las representaciones a nivel de sesión o conexión agrupan paquetes pertenecientes al mismo intercambio, habitualmente utilizando la quíntupla de IP de origen, IP de destino, puerto de origen, puerto de destino y protocolo. Pueden incluir el estado de la conexión, el comportamiento del handshake, metadatos de la aplicación o atributos derivados de TLS/HTTP. Las representaciones a nivel de flujo son más compactas: agregan paquetes en una ventana temporal o registro similar a una conexión, y resumen el comportamiento mediante duración, bytes, paquetes, direccionalidad, tiempos entre llegadas, estadísticas de longitud de paquetes, conteos de flags TCP y periodos activos o inactivos \parencite{Sperotto2010FlowIDS,Umer2017FlowIDS}.

Los registros de estilo NetFlow e IPFIX son el punto de referencia operativo para la monitorización de flujos escalable. Routers, switches, sondas y colectores exportan registros de flujo compactos que pueden retenerse durante largos periodos y procesarse a escala. Estos registros son atractivos para el NIDS porque reducen el volumen y evitan la inspección de la carga útil. También tienen limitaciones porque los despliegues exportan a menudo solo un subconjunto de los campos posibles, pueden muestrear el tráfico y pueden omitir estadísticas más ricas utilizadas en conjuntos de datos académicos. Esto crea una brecha entre la telemetría de flujos en producción y los CSV de benchmark con abundantes características.

Las características al estilo CICFlowMeter se sitúan en el lado más rico y orientado a la investigación de esa brecha. CICIDS2017 proporciona PCAPs y CSV de flujos bidireccionales generados, con características como la duración del flujo, conteos de paquetes y bytes en sentido directo e inverso, estadísticas de tiempos entre llegadas, estadísticas de longitud de paquetes, conteos de flags TCP, agregados de longitud de cabecera, tasas de flujo y temporización activa/inactiva \parencite{Sharafaldin2018CICIDS2017,Lashkari2017CICFlowMeter,CICIDS2017Web}. Estas características son convenientes para ML tabular porque cada flujo se convierte en una fila de anchura fija, pero dependen de la reconstrucción offline de PCAPs y de las decisiones de extracción de características.

La idoneidad de la representación a nivel de flujo es, por tanto, condicional. Es adecuada para esta tesis porque el defensor implementado consume vectores numéricos de anchura fija, porque CICIDS2017 ya está disponible en formato de flujo, y porque la ruta prevista en Fase 2 también extrae CSV de flujos antes de la inferencia. Está limitada porque se pierden la semántica de la carga útil, el ordenamiento exacto de paquetes, el contenido de la aplicación y algunas dependencias temporales de múltiples flujos. Los ataques cuya evidencia es principalmente a nivel de contenido, a nivel de host o de comportamiento a largo plazo pueden ser solo indirectamente visibles para un clasificador de flujos \parencite{Sperotto2010FlowIDS,Umer2017FlowIDS}.

La disponibilidad de características es otra limitación. Una característica calculada por CICFlowMeter a partir de PCAPs completos puede no estar disponible en un exportador NetFlow/IPFIX en producción, o puede calcularse de manera diferente. Sarhan et al. abordan este problema mapeando conjuntos heterogéneos de características de NIDS hacia una representación estándar alineada con NetFlow para mejorar la generalizabilidad y la explicabilidad \parencite{Sarhan2021}. Esta tesis no adopta el conjunto exacto de características estándar de Sarhan, pero sigue la misma motivación metodológica: definir un contrato de características estable, documentar qué está presente o ausente, y evitar tratar un esquema CSV específico de un benchmark como si fuera universal.

El proyecto actual utiliza un esquema canónico de flujo con 76 valores de características y una máscara de presencia/ausencia de 76 valores, produciendo observaciones de 152 dimensiones. La máscara registra si cada valor canónico estaba presente y era válido, en lugar de ocultar silenciosamente los campos ausentes o imputados. Esto importa al pasar de CICIDS2017 a tráfico de laboratorio porque diferentes pipelines de extracción pueden no producir exactamente las mismas columnas. La máscara no hace que los datos faltantes sean inofensivos, y no resuelve el cambio de dominio. Hace auditable la entrada del modelo: un lector puede distinguir un cero medido de un marcador de posición imputado.

La representación a nivel de flujo también modifica el perfil de privacidad y de fugas. La eliminación de las cargas útiles puede reducir la exposición de privacidad, pero los metadatos aún pueden identificar usuarios, servicios, horarios y escenarios. Las direcciones IP, las marcas temporales absolutas, los Flow ID, la identidad de los endpoints, la estructura del día de captura y los puertos usados como proxies de escenario pueden filtrar etiquetas en lugar de representar comportamiento de ataque \parencite{Arp2020DosDontsMLSecurity,Kaufman2011LeakageDataMining}. La tesis trata por tanto la representación y la selección de características como decisiones de diseño relevantes para la seguridad, no solo como detalles de preprocesamiento.

\section{Conjuntos de datos públicos para la detección de intrusiones en red}

Los conjuntos de datos públicos de NIDS permiten la experimentación compartida, la reproducibilidad y la comparación entre trabajos \parencite{Ring2019DatasetSurvey,Thakkar2020DatasetsReview}. También condicionan las conclusiones. Un conjunto de datos refleja su generador de tráfico, su topología, su proceso de captura, su extractor de características, sus etiquetas, sus scripts de ataque y su formato de distribución. El rendimiento sobre un benchmark público es, por tanto, evidencia sobre una tarea controlada, no evidencia directa de preparación para el despliegue \parencite{Apruzzese2022CrossEvaluationNIDS,Layeghy2023CrossDomainNIDS}.

KDDCup99 y NSL-KDD son históricamente importantes porque dieron forma a la evaluación temprana de IDS \parencite{KDDCup1999,Tavallaee2009NSLKDD}. KDDCup99 contiene registros de conexión derivados del tráfico de la era DARPA, con un esquema de 41 características y categorías de ataque como DoS, Probe, R2L y U2R. NSL-KDD redujo algunos problemas de redundancia y facilitó el uso del benchmark, pero mantuvo el mismo origen básico y lógica de características. Estos conjuntos son útiles para comprender la historia de la evaluación de IDS, pero evidencia débil para la defensa moderna basada en flujos: su tráfico, comportamiento de ataques y esquema son antiguos \parencite{Tavallaee2009NSLKDD,Ring2019DatasetSurvey}. En este repositorio, NSL-KDD es contexto histórico, no la ruta final del defensor actual.

UNSW-NB15 es un conjunto de datos de cyber-range más reciente, generado con tráfico benigno y malicioso contemporáneo, capturas sin procesar y características de ingeniería \parencite{MoustafaSlay2015UNSWNB15,Moustafa2016UNSWEval}. Sigue siendo útil porque amplió el campo más allá de los benchmarks al estilo KDD e incluye categorías de ataque más ricas. CICIDS2017 y CSE-CIC-IDS2018 son conjuntos de datos de la familia CIC ampliamente utilizados para la investigación de NIDS basada en flujos; ambos proporcionan escenarios de ataque generados en laboratorio y características de flujo derivadas de capturas de paquetes \parencite{Sharafaldin2018CICIDS2017,CSECICIDS2018AWS}. Bot-IoT y ToN\_IoT amplían el panorama hacia IoT e IIoT, donde el comportamiento de los dispositivos, la telemetría y los patrones de ataque difieren de las redes empresariales \parencite{Koroniotis2019BotIoT,Moustafa2020ToNIoT,Alsaedi2020ToNIoT}.

\begin{table}[htbp]
\centering
\scriptsize
\setlength{\tabcolsep}{3pt}
\begin{tabular}{p{1.85cm}p{1.9cm}p{2.45cm}p{2.65cm}p{2.8cm}}
\toprule
Conjunto de datos & Papel como benchmark & Representación & Principales fortalezas & Principales advertencias \\
\midrule
KDDCup99 & Benchmark histórico heredado & Registros de conexión con 41 características artesanales & Muy ampliamente utilizado; línea base sencilla para la bibliografía clásica de IDS \parencite{KDDCup1999} & Tráfico obsoleto; redundancia; débil como aproximación al NIDS moderno \parencite{Ring2019DatasetSurvey} \\
NSL-KDD & Benchmark heredado depurado & Mismo esquema de conexión al estilo KDD & Reduce algunos problemas de registros duplicados; útil para comparación histórica \parencite{Tavallaee2009NSLKDD} & Sigue heredando las suposiciones antiguas de DARPA/KDD y el diseño de características \\
UNSW-NB15 & Conjunto de datos moderno de laboratorio/cyber-range & PCAPs, registros y características tabulares de ingeniería & Taxonomía de ataques más contemporánea y características más ricas \parencite{MoustafaSlay2015UNSWNB15} & Entorno controlado; no constituye evidencia directa de producción \\
CICIDS2017 & Benchmark interno principal de la tesis & PCAPs y CSV de flujos CICFlowMeter & PCAPs públicos; múltiples días de ataque; ricas características de flujo \parencite{Sharafaldin2018CICIDS2017} & Problemas conocidos de etiquetado, extracción, duplicación y sensibilidad a la partición \parencite{Engelen2021CICIDSIssues,Lanvin2023Errors} \\
CSE-CIC-IDS2018 & Benchmark de mayor tamaño de la familia CIC & PCAPs y flujos al estilo CICFlowMeter & Colaboración más amplia CIC/CSE; comparación moderna útil \parencite{CSECICIDS2018AWS} & Advertencias similares sobre generación en laboratorio y artefactos de flujo \parencite{Liu2022ErrorPrevalenceNIDS} \\
Bot-IoT & Benchmark de IoT/botnets & PCAPs, flujos al estilo Argus, exportaciones CSV & Escenarios de botnet, escaneo, DoS/DDoS y exfiltración orientados a IoT \parencite{Koroniotis2019BotIoT} & Fuerte desbalanceo y artefactos de topología específicos del conjunto de datos \\
ToN\_IoT & Benchmark de telemetría IoT/IIoT & Flujos de red, registros de host y telemetría de sensores & Perspectiva multimodal de IoT/IIoT \parencite{Alsaedi2020ToNIoT} & Entorno de nicho; los artefactos de identificadores y escenarios siguen siendo posibles \\
\bottomrule
\end{tabular}
\caption[Comparación de conjuntos de datos públicos de NIDS]{Comparación compacta de los conjuntos de datos públicos de NIDS relevantes para la tesis.}
\label{tab:nids-datasets}
\end{table}

La Tabla~\ref{tab:nids-datasets} no establece un rango. Muestra que la elección del conjunto de datos debe alinearse con la solidez de las afirmaciones. Los benchmarks históricos sitúan el campo; los conjuntos de datos modernos de laboratorio permiten comparación controlada; los de IoT prueban suposiciones de tráfico diferentes; y el tráfico de laboratorio externo o similar al de producción aporta evidencia de generalización más sólida. Ninguno de estos conjuntos de datos debe tratarse como un sustituto completo de todas las redes reales. Su valor es máximo con preprocesamiento transparente, líneas base sólidas y protocolos que prueben el cambio de distribución, no solo el rendimiento sobre distribuciones aleatorias emparejadas.

\section{CICIDS2017 como benchmark interno principal}

CICIDS2017 es el benchmark interno principal de esta tesis porque proporciona tráfico benigno y de ataque etiquetado, PCAPs sin procesar y CSV de flujos generados con características al estilo CICFlowMeter \parencite{Sharafaldin2018CICIDS2017,Lashkari2017CICFlowMeter,CICIDS2017Web}. Fue diseñado para mejorar conjuntos de datos de IDS anteriores mediante perfiles de tráfico contemporáneos, actividad benigna realista y múltiples escenarios de ataque \parencite{Sharafaldin2018CICIDS2017,Sharafaldin2019Detailed}. Estas propiedades lo hacen adecuado para la experimentación controlada con modelos de NIDS a nivel de flujo.

Su estructura también es atractiva para la tesis. El conjunto abarca múltiples días de captura y familias de ataque, y está disponible tanto como capturas de paquetes sin procesar como en CSV generados. Los PCAPs están más cerca de la evidencia de medición original. Los CSV facilitan el ML porque CICFlowMeter ya reconstruyó los flujos bidireccionales y calculó características estadísticas. Esa comodidad también es un riesgo: al trabajar solo con CSV resumidos, se puede dejar de comprobar cómo se generaron los flujos, cómo se asignaron las etiquetas y si las filas mantienen la independencia entre las particiones de entrenamiento y prueba.

CICFlowMeter reconstruye flujos de trazas de paquetes mediante claves de endpoint, direccionalidad, lógica de timeout y cálculos de características. Pequeñas diferencias en los parámetros de extracción pueden cambiar filas, duraciones, conteos y etiquetas. Engelen et al. muestran problemas prácticos en la reconstrucción de flujos y la consistencia de CICIDS2017 \parencite{Engelen2021CICIDSIssues}. Lanvin et al. reportan errores que afectan al rendimiento de detección y que correcciones aparentemente pequeñas pueden producir diferencias métricas significativas \parencite{Lanvin2023Errors,FaultyUseCICIDS2017}. Liu et al. analizan la prevalencia de errores en CIC-IDS-2017 y CSE-CIC-IDS-2018, mientras Rosay et al. ofrecen un análisis más amplio y una discusión orientada a la reconstrucción \parencite{Liu2022ErrorPrevalenceNIDS,Rosay2022CICIDS2017Analysis}. La crítica de Dube es especialmente relevante para la disciplina en las afirmaciones, porque cuestiona las interpretaciones habituales de puntuaciones elevadas sobre los datos resumidos de CIC-IDS 2017 \parencite{Dube2024FaultyCICIDS2017}.

Estas críticas no hacen inútil a CICIDS2017; hacen más identificable el uso irresponsable. Los patrones comunes incluyen tratar los CSV pregenerados como verdad incuestionable, usar particiones aleatorias por filas de todos los días, conservar identificadores o marcas temporales que exponen el contexto del escenario, informar solo de la exactitud, ignorar el desbalanceo de clases y no documentar si los datos son originales, corregidos, reconstruidos o curados. Las variantes corregidas o reconstruidas pueden ser valiosas, pero no son intercambiables con la distribución original de CSV. Una tesis debe nombrar qué variante utiliza y evitar mezclar resultados de diferentes variantes como si procedieran de un único conjunto de datos.

El repositorio actual responde a estas preocupaciones mediante una copia curada y rastreada, una copia de referencia sin rastrear y limpieza a nivel de cargador en \texttt{src/load\_cicids2017.py}. El adaptador realiza coerción numérica, tratamiento de valores no válidos, mapeo canónico y exclusión anti-fugas. También ajusta el escalado sobre la partición de entrenamiento y aplica la transformación ajustada a prueba, en vez de ajustar sobre todos los datos. El mapeo canónico en \texttt{src/canonical\_schema.py} excluye IPs, marcas temporales, Flow ID y puertos específicos de las características canónicas.

El código de evaluación admite particiones aleatorias, por CSV/día y validación \textit{leave-one-CSV-out}. El rendimiento con partición aleatoria sirve para comprobar que modelo e implementación aprenden el benchmark interno. Las particiones por día/archivo más difíciles, las comprobaciones de etiquetas barajadas, la validación \textit{leave-one-CSV-out} y la inferencia offline sobre flujos de laboratorio prueban modos de fallo distintos. Una puntuación alta en el escenario más sencillo no basta para una afirmación de despliegue. En esta tesis, CICIDS2017 debe leerse como benchmark interno principal, no como sustituto de la validación con tráfico capturado por el operador en un laboratorio doméstico cerrado y no adversarial.

\section{Aprendizaje supervisado y aprendizaje profundo para NIDS}

La mayor parte del NIDS basado en flujos se formula como aprendizaje supervisado: cada registro de tráfico se mapea a una etiqueta benigna, maliciosa o de familia de ataque \parencite{LiuLang2019IDSSurvey,Ahmad2021}. Los métodos clásicos incluyen árboles de decisión, Random Forests, \textit{SVMs}, \textit{k-NN}, regresión logística, \textit{Naive Bayes} y conjuntos de árboles con \textit{gradient boosting}. Difieren en sus suposiciones, coste operativo y requisitos de memoria, entrenamiento e inferencia. La regresión logística es simple e interpretable y ofrece una referencia lineal, pero está limitada ante interacciones no lineales. \textit{Naive Bayes} es rápido y sirve como línea base débil, aunque sus supuestos de independencia son irrealistas para estadísticas de flujo correlacionadas. k-NN puede capturar vecindades locales, pero escala mal con tablas grandes y es sensible al escalado de características. Las SVMs pueden ser potentes en conjuntos pequeños, pero resultan costosas con tablas de flujos grandes y exigen escalado cuidadoso de características y selección del kernel.

Los modelos basados en árboles son especialmente importantes para NIDS tabular. Los árboles de decisión y Random Forests manejan interacciones no lineales, escalas heterogéneas, variables irrelevantes y distribuciones mixtas \parencite{LiuLang2019IDSSurvey,Apruzzese2022CrossEvaluationNIDS}. Las familias con \textit{gradient boosting}, como \textit{XGBoost}, \textit{LightGBM} y \textit{CatBoost}, son a menudo aprendices tabulares potentes, pero este capítulo no necesita afirmaciones separadas sobre cada implementación si la tesis no las evalúa. El punto metodológico importante es que un clasificador de flujos basado en RL debe compararse con líneas base supervisadas tabulares competentes, no solo con métodos débiles u obsoletos.

Esto es directamente relevante para la tesis. Si cada observación es una fila de flujo etiquetada y la recompensa deriva de la corrección de la clasificación, un clasificador supervisado es la línea base natural. Un agente QRDQN no puede evaluarse únicamente contra su curva de recompensa, pérdida de entrenamiento o una línea base de clase mayoritaria; debe compararse con métodos supervisados sólidos bajo el mismo esquema de características, protocolo de partición, preprocesamiento y métricas. El repositorio ya incluye una línea base de Random Forest alineada con las particiones de QRDQN y artefactos por barrido. La comparación debe limitarse a métricas respaldadas por esos artefactos y al mismo protocolo, en vez de inferir superioridad general.

El aprendizaje profundo amplía el espacio de diseño mediante MLPs, CNNs, RNNs, LSTMs, autoencoders, mecanismos de atención, modelos al estilo transformer y enfoques basados en grafos \parencite{Xiao2021DLIDSSurvey,Asadullah2023MLIDSSurvey,Sayeeduddin2022AINIDS,Nguyen2022BERTFlowNIDS}. Los MLPs tratan las características de flujo como vectores numéricos genéricos; las CNNs se utilizan cuando bytes de paquetes, matrices de características o patrones locales se codifican en forma estructurada. Las RNNs y LSTMs se orientan a secuencias temporales de paquetes, flujos o eventos. Los autoencoders apoyan frecuentemente la detección de anomalías al aprender representaciones comprimidas del comportamiento benigno y señalar errores de reconstrucción. Los modelos al estilo transformer y las redes neuronales en grafo son direcciones más recientes para representaciones de tráfico secuenciales y estructurales.

Estas arquitecturas pueden alcanzar puntuaciones de benchmark muy elevadas, sobre todo en conjuntos públicos, pero las revisiones identifican debilidades recurrentes: preprocesamiento poco claro, particiones aleatorias favorables, análisis insuficiente del desbalanceo, validación externa limitada y reproducibilidad insuficiente \parencite{LiuLang2019IDSSurvey,Ring2019DatasetSurvey,Arp2020DosDontsMLSecurity}. Un modelo profundo puede sobreajustarse a los artefactos del benchmark igual que un modelo superficial. La familia de modelos por sí sola no resuelve fugas, desbalanceo de clases, disponibilidad de características ni cambio de dominio.

La estandarización de características es otra lección importante del aprendizaje supervisado. Diferentes conjuntos de datos y herramientas exponen conjuntos solapados, pero no idénticos, lo que dificulta la comparación entre conjuntos. Sarhan et al. propusieron un mapeo de características alineado con NetFlow para mejorar la generalizabilidad y la explicabilidad en NIDS \parencite{Sarhan2021}. La tesis no adopta ese estándar exacto, pero sí su motivación: un esquema canónico fijo hace auditable la interfaz del modelo y permite comparar las entradas de flujo de CICIDS2017 y de laboratorio a través de un único contrato de preprocesamiento. La máscara de presencia/ausencia extiende esta idea al hacer visibles las características ausentes al modelo y al evaluador.

\section{Fundamentos del aprendizaje por refuerzo}

El aprendizaje por refuerzo estudia cómo un agente selecciona acciones en un entorno para maximizar el retorno \parencite{SuttonBarto2018RL}. Sus conceptos centrales son los estados u observaciones, las acciones, las recompensas, las políticas, las funciones de valor y las transiciones. En un proceso de decisión de Markov, el estado actual resume la información necesaria para la toma de decisiones futura, y las acciones influyen tanto en la recompensa inmediata como en los estados posteriores \parencite{SuttonBarto2018RL}. Esta propiedad secuencial es lo que distingue el RL de la clasificación supervisada ordinaria.

Una política especifica cómo el agente elige acciones a partir de las observaciones. Una función de valor estima el retorno esperado desde un estado, mientras que una función de valor-acción estima el retorno esperado tras tomar una acción particular en un estado. Las recompensas definen el objetivo de la tarea, y el retorno agrega las recompensas futuras, generalmente con descuento. Estas definiciones importan porque cambiar la función de recompensa cambia aquello para lo que el agente está optimizado. En un contexto de seguridad, una recompensa que penaliza más los falsos negativos que los falsos positivos codifica una preferencia operativa concreta; no es un hecho neutral sobre todas las redes.

La distinción entre tareas episódicas y continuas también es relevante. Las tareas episódicas tienen límites y estados terminales, como un episodio de juego o un escenario simulado de respuesta a incidentes. Las tareas continuas modelan la operación continua, como la monitorización continua de la red. Una formulación estática de conjunto de datos como entorno impone a menudo episodios artificiales: un episodio puede ser un número fijo de filas muestreadas, un recorrido por un conjunto de datos o un archivo. Eso es aceptable para la experimentación, pero es más débil que un entorno donde la acción del defensor modifica el estado siguiente.

El Q-learning clásico es un método de diferencia temporal model-free y off-policy que aprende valores de acción sin requerir un modelo de la dinámica del entorno \parencite{WatkinsDayan1992QLearning}. Los métodos model-free aprenden directamente de la interacción o de transiciones registradas, mientras que los métodos basados en modelo aprenden o utilizan un modelo de transición para la planificación. Los métodos on-policy aprenden sobre la política que está generando el comportamiento actual; los métodos off-policy pueden aprender sobre una política objetivo mientras el comportamiento es generado de manera diferente. DQN y QRDQN heredan el enfoque basado en valores y off-policy. La repetición de experiencias, introducida anteriormente como mecanismo para reutilizar experiencias pasadas, también se volvió central en los sistemas de Q-learning profundo posteriores \parencite{Lin1992ReactiveAgents}.

El RL tabular no es apropiado cuando las observaciones son vectores numéricos de alta dimensión. Una tabla de búsqueda sobre todas las posibles combinaciones de características de flujo es inviable, y pequeños cambios numéricos podrían crear efectivamente estados nuevos. La aproximación de funciones es, por tanto, necesaria. El DRL utiliza redes neuronales para aproximar políticas o funciones de valor sobre espacios de observación grandes. En esta tesis, la observación es un vector de 152 dimensiones: 76 valores canónicos de flujo más 76 valores de máscara de presencia/ausencia. Eso justifica un algoritmo de la familia DQN más que un algoritmo de RL tabular, pero no justifica por sí solo el RL frente al aprendizaje supervisado.

La ciberseguridad puede contener problemas de RL secuenciales genuinos. Un defensor puede aislar un host, parchear un servicio, cambiar el enrutamiento, desplegar una contramedida o elegir dónde inspeccionar a continuación, y esas acciones pueden afectar a las observaciones posteriores y a las oportunidades del atacante. Sin embargo, un conjunto de datos estático de flujos etiquetados es un escenario más débil. Si se muestrea una fila de flujo, se clasifica y va seguida de otra fila que no es causalmente afectada por la acción anterior, el problema se parece más a una clasificación sensible al coste o a un proceso de decisión contextual que a una defensa cibernética autónoma completa \parencite{Levine2020OfflineRL,Fujimoto2019BCQ}. La tesis debe mantener esta distinción de forma explícita.

\section{Q-Learning profundo y aprendizaje por refuerzo distribucional}

El Q-Learning Profundo sustituye una función de valor-acción tabular por una red neuronal, haciendo aplicable el RL basado en valores a observaciones de alta dimensión \parencite{Mnih2015DQN}. DQN popularizó dos mecanismos estabilizadores: la memoria de repetición y una red objetivo separada \parencite{Mnih2015DQN}. La memoria de repetición rompe las correlaciones a corto plazo muestreando transiciones pasadas para el entrenamiento y mejora la reutilización de datos. Una red objetivo estabiliza los objetivos de bootstrapping manteniendo el cálculo del objetivo separado de los pesos de la red online, que cambian rápidamente.

Variantes posteriores de DQN abordaron debilidades conocidas. Double DQN reduce el sesgo de sobreestimación desacoplando la selección de acciones de la evaluación del objetivo \parencite{VanHasselt2016DoubleDQN}. Dueling DQN separa la estimación del valor del estado y la ventaja de la acción, lo que puede ayudar cuando muchas acciones tienen valores similares en un estado \parencite{Wang2016DuelingDQN}. Prioritized Experience Replay muestrea transiciones según la señal de aprendizaje en lugar de hacerlo de forma uniforme \parencite{Schaul2016PER}. Rainbow combinó varias mejoras en un único agente de la familia DQN, incluyendo double Q-learning, redes dueling, repetición priorizada, retornos de múltiples pasos, aprendizaje distribucional y exploración con ruido \parencite{Hessel2018Rainbow}.

El aprendizaje por refuerzo distribucional cambia el objeto que se estima. En lugar de estimar únicamente el retorno esperado, aproxima la distribución de los retornos posibles \parencite{Bellemare2017DistributionalRL}. C51 representa la distribución del retorno sobre un soporte categórico fijo. QRDQN usa regresión cuantílica para aproximar los cuantiles del retorno, evitando la misma suposición de soporte fijo \parencite{Dabney2018QRDQN}. QRDQN no es, por tanto, un paradigma separado de DQN; es un miembro distribucional y basado en valores de la familia DQN.

La motivación de seguridad es intuitiva, pero debe formularse con cuidado. Los falsos positivos y los falsos negativos tienen consecuencias asimétricas, y los errores raros de alto coste pueden importar más de lo que sugiere una puntuación media. Una estimación distribucional del valor puede ser útil para estudiar la estructura del retorno bajo recompensas asimétricas. Puede revelar si una decisión tiene una distribución de retornos estrecha o amplia bajo el modelo aprendido. Sin embargo, el RL distribucional no proporciona automáticamente incertidumbre calibrada, seguridad operativa, comportamiento sensible al riesgo ni superioridad sobre el DQN escalar. Esas propiedades requerirían reglas de decisión explícitas, calibración, objetivos de riesgo en la cola o criterios de evaluación más allá de simplemente entrenar un agente distribucional \parencite{Bellemare2017DistributionalRL,Dabney2018QRDQN}.

En esta tesis, QRDQN es una elección algorítmica exploratoria. Es adecuado para un espacio de acciones binario discreto y observaciones de flujo de dimensión moderadamente alta, y está disponible a través de la implementación de SB3-Contrib utilizada por el proyecto \parencite{SB3ContribQRDQNDocs,GymnasiumDocs}. La elección algorítmica es coherente: dos acciones, observaciones numéricas, aprendizaje de valores off-policy, repetición y un objetivo distribucional. La afirmación evidencial sigue siendo empírica. QRDQN no debe describirse como probadamente superior para NIDS a menos que las comparaciones con el mismo protocolo frente a DQN, Random Forest y otras líneas base apoyen esa afirmación.

\section{Aprendizaje por refuerzo para la detección de intrusiones y la defensa cibernética}

El RL y el DRL para la detección de intrusiones ya forman una bibliografía activa \parencite{Yang2024DRLNIDSSurvey,Kheddar2025RLIDSReview,Gueriani2024DRLIoTIDSSurvey,Adawadkar2022CyberSecurityRLSurvey}. Los trabajos previos incluyen clasificadores de intrusiones al estilo DQN, enfoques actor-crítico, formulaciones de entorno adversarial, agentes de selección de características, sistemas orientados a SDN o IoT, y simulaciones de defensa cibernética autónoma. El campo es heterogéneo, por lo que los trabajos no deben tratarse como directamente comparables a menos que su formulación de tarea, conjuntos de datos, acciones, recompensas y protocolos de validación estén alineados.

Una rama está próxima a esta tesis: la detección de intrusiones con conjunto de datos como entorno. Lopez-Martin et al. reformularon la detección supervisada de intrusiones como un problema de RL profundo tratando los registros etiquetados como estados, las predicciones de clase como acciones y la recompensa como función de la corrección \parencite{LopezMartin2020DRLIDS}. Su trabajo posterior extendió una formulación relacionada de RL offline a varios conjuntos de datos, incluyendo CICIDS2017 y UNSW-NB15 \parencite{LopezMartin2021RBFOfflineRL}. Ren et al. utilizaron DRL para la selección de características y la clasificación en CSE-CIC-IDS2018 \parencite{Ren2022IDRDRL}. Estos trabajos muestran que las formulaciones de RL para IDS tienen precedentes.

Otros trabajos de RL orientados a IDS exploran DQN, Double DQN, DQN adversarial, actor-crítico, SAC, DDPG y diseños al estilo Rainbow \parencite{Caminero2019AERLIDS,AEDDQN2020NSLKDD,AESAC2023RLIDS,Alavizadeh2021,Hossain2025DQIDS}. Los detalles varían ampliamente. Algunos trabajos utilizan clasificación binaria, otros etiquetas de ataque multiclase, selección de características, control de umbral o acciones específicas de SDN. Muchos siguen evaluándose sobre conjuntos de datos públicos estáticos. Eso los convierte en antecedentes relevantes para la tesis, pero también limita la comparabilidad directa. Un trabajo que usa NSL-KDD con particiones aleatorias y recompensas de corrección no es evidencia de que un defensor QRDQN offline generalice de CICIDS2017 a tráfico de laboratorio.

Una segunda rama estudia la defensa cibernética genuinamente secuencial. Los enfoques POMDP y de grafo de ataque modelan acciones defensivas bajo incertidumbre \parencite{Hu2020POMDPDefense}. CybORG proporciona un entorno al estilo gym para agentes cibernéticos autónomos \parencite{Standen2021CybORG}. Los trabajos más amplios de defensa cibernética autónoma discuten agentes que seleccionan contramedidas en redes simuladas o emuladas \parencite{Palmer2023_weak,CSET2023AutonomousCyberDefense}. Estos escenarios son conceptualmente diferentes de esta tesis porque las acciones del defensor pueden alterar el estado futuro. Representan una forma más fuerte de problema de RL que la clasificación estática por filas.

\begingroup
\scriptsize
\setlength{\tabcolsep}{2pt}
\renewcommand{\arraystretch}{1.08}
\begin{longtable}{@{}p{1.75cm}p{2.05cm}p{1.95cm}p{3.55cm}p{3.00cm}@{}}
\caption[Trabajos de RL para NIDS]{Síntesis de trabajos de RL/DRL para NIDS y límites de comparabilidad con esta tesis. La relación indicada se infiere de cada fuente y de la formulación QRDQN offline sobre flujos de CICIDS2017 adoptada en este trabajo.}\label{tab:rl-nids-related-work}\\
\toprule
\textbf{Fuente} & \textbf{Formulación de RL} & \textbf{Datos / ámbito} & \textbf{Límite de evaluación} & \textbf{Relación / diferencia con esta tesis} \\
\midrule
\endfirsthead
\multicolumn{5}{@{}l}{\small\itshape Tabla \thetable{} (continuación)} \\
\toprule
\textbf{Fuente} & \textbf{Formulación de RL} & \textbf{Datos / ámbito} & \textbf{Límite de evaluación} & \textbf{Relación / diferencia con esta tesis} \\
\midrule
\endhead
\midrule
\multicolumn{5}{r@{}}{\small\itshape Continúa en la página siguiente} \\
\endfoot
\bottomrule
\endlastfoot
\textcite{NIDSRL2023} & RL; el algoritmo no es verificable desde el registro primario público. & NIDS genérico; datos no verificables. & El registro no expone protocolo ni métricas comparables. & Antecedente temático, no comparación directa con QRDQN/CICIDS2017. \\
\textcite{RLTechniques2023NIDS} & Técnicas de RL; algoritmo no verificable desde el registro primario público. & NIDS genérico; datos no verificables. & El registro no expone protocolo ni métricas comparables. & Antecedente genérico, sin comparabilidad empírica directa. \\
\textcite{Sanusi2023DRLIDS} & DRL para IDS binario. & NSL-KDD; la tarea multiclase se reduce a binaria. & Tesis offline sobre benchmark; no evidencia de despliegue. & Antecedente de IDS binario con DRL, pero datos e implementación distintos. \\
\textcite{Umer2022RLRLIDS} & DQN y gradiente de política con entrenamiento adversarial. & IDS a nivel de paquete y flujo; CICDDoS2019. & Evaluación de benchmark; sin evidencia de despliegue operativo. & Diseño paquete/flujo más rico y datos distintos de QRDQN/CICIDS2017. \\
\textcite{Cevallos2023DRLIDSBP} & Revisión/taxonomía de DRL-IDS, no un algoritmo único. & Literatura de IDS en IoT. & Síntesis bibliográfica, sin benchmark propio. & Contexto de decisiones DRL en IoT, no línea base empírica. \\
\textcite{DDPG2025AttackDetection} & DDPG. & Detección de ataques de red; datos no verificables. & El registro no expone protocolo ni métricas comparables. & Familia de RL distinta; solo antecedente a nivel de título. \\
\textcite{DRLIDSSDN2025} & DRL-IDS con LFTS-RNN y PC-JTFOA. & SDN; NSL-KDD y WPPD. & Particiones de benchmark, no despliegue SDN vivo. & Tarea, modelo y datos distintos; no equivale a una evaluación QRDQN offline. \\
\textcite{HCLRIDS2025IoMT} & CNN-LSTM con DQN y PPO. & IoMT sanitario; CICIoMT2024. & Benchmark especializado; el trabajo deja la detección \textit{zero-day} como trabajo futuro. & Propuesta IoMT específica, no línea base comparable con CICIDS2017. \\
\end{longtable}
\endgroup

La tesis pertenece, por tanto, principalmente a la primera rama: detección offline a nivel de flujo formulada a través de una interfaz de RL. Puede referenciar la defensa cibernética autónoma como contexto más amplio y dirección futura, pero no debe implicar que el sistema actual realice aprendizaje multiagente, adversarial o de control de red en vivo. La formulación más precisa es que la tesis implementa un benchmark de NIDS con conjunto de datos como entorno y una interfaz de acción de seguridad binaria sensible al coste, no un defensor de cyber-range completo.

\section{Clasificación como RL y decisiones de seguridad binarias}

La formulación \texttt{PERMIT}/\texttt{BLOCK} pertenece a la discusión de clasificación como RL. En el entorno actual, cada observación de flujo se presenta al agente, la acción es binaria y la recompensa se calcula a partir de la acción y la etiqueta de verdad. \texttt{PERMIT} corresponde a aceptar el tráfico como benigno, mientras que \texttt{BLOCK} corresponde a identificarlo como ataque. Un falso negativo permite un flujo de ataque; un falso positivo bloquea tráfico benigno.

Este enfoque tiene un beneficio metodológico real: hace explícitas las decisiones de seguridad sensibles al coste. Los errores de IDS no son simétricos en la práctica. Omitir un ataque y bloquear tráfico benigno tienen significados operativos diferentes, y la compensación preferida depende del entorno defendido \parencite{Lee2002CostSensitiveIDS,Cardenas2006IDSFramework}. El diseño de la recompensa puede codificar esta asimetría directamente, lo que es útil para una tesis que estudia reglas de decisión defensivas en lugar de solo la predicción de etiquetas.

La limitación es igualmente importante. Si el conjunto de datos no reacciona a la acción, entonces \texttt{PERMIT}/\texttt{BLOCK} no es una acción de control activa. Es una etiqueta de decisión offline sobre una fila de flujo registrada o extraída. La fila siguiente no cambia porque el agente haya bloqueado o permitido la fila anterior. La formulación es, por tanto, más próxima a una clasificación con recompensa diseñada que a un firewall, un IPS, un controlador SDN o un agente de defensa cibernética autónoma. Esto no es un defecto si se declara con claridad. Solo se convierte en un defecto si el capítulo reclama en exceso el bloqueo en vivo, la preparación para el despliegue o el control secuencial completo.

La función de recompensa también necesita una interpretación cuidadosa. Una recompensa positiva para un verdadero positivo y una penalización mayor para un falso negativo puede expresar la idea de que los ataques omitidos son costosos. Una penalización por falso positivo puede expresar el coste de interrumpir el tráfico benigno o de sobrecargar a los analistas. Una recompensa de cero o pequeño valor para los verdaderos negativos puede mantener el foco en la detección de ataques. Estas son decisiones de diseño, no economía universal de seguridad. Diferentes organizaciones elegirían distintos costes dependiendo de la continuidad del negocio, el modelo de amenaza, la capacidad del analista y la tolerancia a los ataques omitidos.

Dado que la formulación está próxima a la clasificación sensible al coste, las líneas base supervisadas son obligatorias. Un Random Forest sólido, e idealmente otras líneas base tabulares, ayudan a determinar si QRDQN añade valor empírico más allá de las características canónicas, el preprocesamiento y el diseño de la recompensa. Las mismas métricas deben calcularse tanto a partir de las salidas de RL como de las supervisadas. La recompensa de entrenamiento sola no es suficiente, porque una curva de recompensa puede mejorar mientras la precisión, el recall o el comportamiento de falsos positivos sigue siendo operativamente deficiente. La tesis debe interpretar cualquier ventaja de RL como específica del benchmark y el protocolo, a menos que se disponga de evidencia más amplia.

\section{Riesgos metodológicos en NIDS basado en ML/DL/RL}

La bibliografía de NIDS advierte repetidamente contra sobreinterpretar la exactitud de benchmark. Sommer y Paxson señalan que el ML para detección de intrusiones es difícil de validar en entornos realistas: el tráfico benigno es diverso, los ataques raros, las etiquetas difíciles de obtener y las tasas base operativas difieren de las del benchmark \parencite{SommerPaxson2010ClosedWorld}. Arp et al. identifican riesgos afines en ML para seguridad, como fuga de información, sesgo de muestreo, métricas inapropiadas y diseño experimental débil \parencite{Arp2020DosDontsMLSecurity}. La evidencia sobre generalización entre conjuntos y sobre la elección de datos de entrenamiento refuerza que un resultado dentro de un benchmark no basta para inferir el comportamiento en otro entorno \parencite{Layeghy2023CrossDomainNIDS}.

La fuga de información abarca más que poner por error la misma fila en entrenamiento y prueba: incluye cualquier información relacionada con el objetivo que no estaría legítimamente disponible al predecir \parencite{Kaufman2011LeakageDataMining}. En NIDS puede entrar por identificadores de endpoint, marcas temporales absolutas, Flow ID, puertos específicos del escenario, preprocesamiento global ajustado antes de particionar, flujos duplicados o casi duplicados, y particiones que comparten contexto de captura. Un modelo que aprende que una IP, puerto, día o Flow ID concreto corresponde a un ataque puede puntuar bien sin aprender comportamiento malicioso portable.

La fuga temporal y de escenario es especialmente importante. Los conjuntos públicos de NIDS suelen organizarse por día, archivo de captura, ventana de ataque o escenario con script. Una partición aleatoria por filas puede dejar tráfico muy relacionado a ambos lados de la frontera de entrenamiento y prueba. Aunque las filas no sean duplicados exactos, pueden compartir hosts, herramientas de ataque, regímenes de temporización o artefactos de extracción. Esto importa en CICIDS2017 porque trabajos posteriores reportan errores o artefactos en extracción de flujos, etiquetado y uso de CSV resumidos \parencite{Engelen2021CICIDSIssues,Lanvin2023Errors,Liu2022ErrorPrevalenceNIDS,Rosay2022CICIDS2017Analysis,Dube2024FaultyCICIDS2017}.

El preprocesamiento también puede introducir fugas. El escalado, la imputación, la selección de características, el recorte de valores atípicos o la reducción de dimensionalidad deben ajustarse sobre entrenamiento y aplicarse después a los datos de validación o de prueba. Si un escalador se ajusta sobre todos los datos antes de particionar, la información distribucional de prueba entra en el entrenamiento. Puede parecer menos grave que filtrar etiquetas, pero con estructura fuerte por día o escenario también puede inflar la confianza. Por ello, los objetos \texttt{StandardScaler} ajustados sobre entrenamiento en el repositorio forman parte del diseño metodológico, no solo de la codificación.

El desbalanceo de clases y el problema de la tasa base complican la interpretación. En redes operativas, los ataques suelen ser raros frente al tráfico benigno. Un detector puede lograr exactitud elevada mientras omite muchos ataques, o alto recall mientras genera demasiadas alertas falsas para investigarlas. El argumento de la tasa base de Axelsson sigue siendo relevante: incluso tasas bajas de falsos positivos pueden dominar los flujos de alertas cuando la prevalencia real de ataque es baja \parencite{Axelsson1999BaseRate}. Por eso la exactitud no debe ser la métrica principal para NIDS.

El RL introduce riesgos adicionales: el entrenamiento profundo puede ser estocástico y sensible a semillas aleatorias, hiperparámetros, modelado de la recompensa, configuración del buffer de repetición, exploración y detalles del entorno. En un conjunto de datos estático usado como entorno, buenos resultados pueden reflejar las mismas ventajas de características y partición que benefician a modelos supervisados. Una afirmación de superioridad de QRDQN requiere el mismo protocolo de comparación y reconocer la varianza entre ejecuciones cuando no hay semillas repetidas. Las ganancias de una única ejecución deben tratarse como preliminares salvo que las sostengan experimentos repetidos o artefactos de validación sólidos.

\section{Protocolos de evaluación, métricas y reproducibilidad}

La evaluación debe ajustarse a la solidez de la afirmación. Las particiones aleatorias sobre el mismo conjunto de datos son aceptables para comprobaciones internas de aprendizaje, pero son evidencia débil para el despliegue. Una escalera de validación práctica para esta tesis comienza con particiones aleatorias estratificadas, avanza hacia particiones temporales o por CSV/día, la validación leave-one-scenario o leave-one-CSV-out, las pruebas entre conjuntos de datos con un esquema de características armonizado, y finalmente el tráfico capturado de forma independiente o similar al de producción. La tesis no alcanza este último peldaño; su inferencia sobre tráfico de laboratorio, generado por el operador en un entorno doméstico cerrado y no adversarial, se sitúa por debajo. Cada paso aumenta la separación distribucional entre datos de entrenamiento y prueba, y por tanto permite afirmaciones de generalización más sólidas.

El repositorio actual refleja esta escalera. El entrenamiento con partición aleatoria verifica que la implementación de QRDQN, el esquema de características, la función de recompensa y el bucle de entrenamiento pueden aprender el benchmark controlado. La comprobación B con etiquetas barajadas comprueba si el rendimiento colapsa hacia una línea base de clase mayoritaria cuando las etiquetas ya no contienen señal real. La comprobación C utiliza una partición por CSV/día más difícil. El script leave-one-CSV-out comprueba si retener un CSV original de CICIDS2017 modifica el comportamiento. La inferencia offline de Fase 2 sobre CSV de flujos de laboratorio prueba la cuestión separada de si un modelo entrenado sobre el benchmark puede procesar flujos capturados por el operador en un laboratorio doméstico cerrado.

Las métricas también deben interpretarse operativamente. La exactitud es útil pero insuficiente para NIDS desequilibrados. La precisión indica cuán fiables son las alertas positivas. El recall indica qué proporción del tráfico de ataque es detectada. F1 proporciona un equilibrio compacto, pero oculta la compensación real entre falsos positivos y falsos negativos. La tasa de falsos positivos mapea al tráfico benigno bloqueado o alertado incorrectamente. La tasa de falsos negativos mapea al tráfico malicioso omitido por el detector \parencite{Axelsson1999BaseRate,Fawcett2006ROC,Saito2015PRROC}. Las matrices de confusión, la precisión, el recall, F1, TP, FP, FN, TN, FPR y FNR por clase deben preferirse sobre un único número resumen \parencite{Milenkoski2015IDSEvaluationSurvey}.

La evaluación sensible al coste es relevante porque el punto de operación correcto depende del contexto. La guía del NIST y la bibliografía de evaluación de IDS enfatizan ambas las compensaciones entre falsos positivos y falsos negativos \parencite{ScarfoneMell2007,Cardenas2006IDSFramework}. En esta tesis, el diseño de recompensa asimétrico debe tratarse como una suposición de escenario y evaluarse empíricamente, no como un modelo de coste universal para todas las redes. Si la recompensa penaliza fuertemente los falsos negativos, la política resultante puede preferir bloquear más flujos. Ese comportamiento debe juzgarse mediante precisión, recall, tasa de falsos positivos y las suposiciones de coste específicas declaradas en el experimento.

La reproducibilidad es parte de la evidencia, no un detalle administrativo. Los resultados de ML en ciberseguridad deben preservar la versión del conjunto de datos, la variante exacta del conjunto de datos, el esquema de características, los campos excluidos, el orden del preprocesamiento, la lógica de partición, el estado del escalador, las semillas aleatorias, la configuración del modelo, la configuración de la recompensa, el identificador de ejecución, las métricas y las predicciones donde sea factible. Olszewski et al. muestran que la reproducibilidad sigue siendo un problema importante en la investigación de ML en seguridad \parencite{Olszewski2023ReproSecurityML}. El uso en el repositorio de configuraciones guardadas, métricas, artefactos del modelo, salidas de validación y carpetas de ejecución es, por tanto, una fortaleza metodológica, pero los artefactos faltantes deben seguir claramente marcados como pendientes.

Esta disciplina en los informes también protege contra las afirmaciones excesivas accidentales. Un resultado vinculado a la ejecución C03 de QRDQN con partición aleatoria completa es un resultado para esa ejecución, con esa partición, configuración de recompensa, ruta de preprocesamiento y estado de artefactos. No debe convertirse silenciosamente en una afirmación sobre todas las configuraciones de QRDQN, todas las variantes de CICIDS2017 o todo el tráfico de red. Cuanto más sólida sea la afirmación, más completo debe ser el rastro de artefactos.
% Esto de arriba sobre artefactos faltantes habrá que quitarlo, nunca habrá artefactos faltantes, ya que si faltan, es que no están dentro del alcance de la tesis, y no hay que decirlo. La tesis debe ser clara sobre lo que está dentro del alcance y lo que no. De todos modos, esos artefactos se generarán pronto por el usuario ejecutando los nuevos experimentos.

\section{Eficiencia de datos y evaluación a escala de entrenamiento}

La eficiencia de datos pregunta cuánto tráfico etiquetado se necesita para alcanzar un rendimiento útil, y cómo cambia la curva de aprendizaje a medida que crece el conjunto de entrenamiento. Esta pregunta es importante para NIDS porque el tráfico de ataque etiquetado es costoso, las distribuciones de tráfico cambian y la abundancia de benchmarks públicos no implica abundancia de datos de despliegue \parencite{Ring2019DatasetSurvey,Apruzzese2022CrossEvaluationNIDS}. También es importante para el RL porque los métodos de RL profundo pueden ser exigentes en muestras, especialmente cuando las señales de recompensa son escasas, ruidosas, retardadas o altamente diseñadas \parencite{SuttonBarto2018RL,Hessel2018Rainbow}.

La bibliografía existente de RL-IDS con frecuencia entrena sobre conjuntos de datos públicos completos e informa de puntuaciones finales en el benchmark, mientras que las ablaciones controladas a escala de entrenamiento son menos comunes \parencite{Yang2024DRLNIDSSurvey,Kheddar2025RLIDSReview}. El trabajo de NIDS supervisado con pocos ejemplos y aprendizaje incremental por clases aborda preguntas relacionadas desde una perspectiva diferente \parencite{DiMonda2024FewShotNIDS}, pero la pregunta de la tesis es más estrecha: bajo un mismo esquema de características y protocolo de evaluación, ¿cómo escala QRDQN con las filas de entrenamiento disponibles, y cómo se compara eso con las líneas base supervisadas?

La respuesta importa porque la elección del modelo y las necesidades de datos están conectadas. Si un Random Forest alcanza un rendimiento estable con muchas menos filas que QRDQN bajo la misma partición, eso es evidencia importante incluso si QRDQN lo alcanza posteriormente. Si QRDQN es más sensible a los datos reducidos o al desbalanceo de la clase de ataque, la tesis debe reportarlo como una limitación. Si QRDQN es competitivo con presupuestos menores, la evidencia debe seguir vinculada al benchmark y al protocolo de validación, no generalizada a todos los escenarios de NIDS.

La tesis posiciona los experimentos a escala de entrenamiento como evidencia metodológica. Presupuestos como 100k, 250k, 500k, 1M y 2M muestras pueden comprobar si la formulación de RL requiere sustancialmente más datos que una línea base supervisada para alcanzar un rendimiento comparable. Tales experimentos no deben enmarcarse como prueba de aprendizaje con pocos ejemplos ni de preparación para el despliegue. Son evidencia interna sobre los requisitos de datos de esta formulación particular. La salida útil es una comparación de curvas de aprendizaje con preprocesamiento consistente, semillas aleatorias donde sea factible, métricas compartidas y rutas de artefactos explícitas.

\section{Brecha de investigación y posicionamiento de esta tesis}

La bibliografía respalda una brecha de investigación estrecha pero significativa. El NIDS y el ML basado en flujos son áreas maduras; conjuntos de datos públicos como CICIDS2017 son útiles pero frágiles; las líneas base tabulares supervisadas son sólidas; el RL y el DRL para IDS ya existen; y la defensa cibernética autónoma es un campo más amplio y más secuencial que la implementación actual \parencite{Ring2019DatasetSurvey,LiuLang2019IDSSurvey,Yang2024DRLNIDSSurvey,Kheddar2025RLIDSReview}. La brecha no es la ausencia de RL para la detección de intrusiones.

La contribución de esta tesis es un estudio experimental reproducible de una formulación offline de decisión de seguridad binaria a nivel de flujo. El sistema mapea CICIDS2017 y entradas posteriores de flujos de laboratorio en un esquema canónico fijo de 76 características, aumenta las observaciones con una máscara de presencia/ausencia para producir entradas de 152 dimensiones, expone las filas a través de un entorno compatible con Gymnasium, y entrena un agente QRDQN para elegir entre \texttt{PERMIT} y \texttt{BLOCK}. Esto hace explícita la interfaz de decisión sensible al coste mientras preserva una obligación de comparación clara contra las líneas base supervisadas.

La tesis también está posicionada metodológicamente. Trata a CICIDS2017 como el benchmark interno principal, no como prueba de preparación para la producción. Separa los resultados con partición aleatoria de la validación más estricta. Utiliza preprocesamiento con control de fugas y comprobaciones de validación para reducir el aprendizaje evidente de artefactos. Trata el tráfico de laboratorio capturado externamente como el siguiente paso evidencial preferido, reconociendo al mismo tiempo que la inferencia actual de Fase 2 es offline y no realiza bloqueo activo.

La brecha final es, por tanto, una brecha de integración: un benchmark QRDQN reproducible, con control de fugas y sensible al coste para el NIDS a nivel de flujo, evaluado contra líneas base supervisadas e interpretado a través de una escalera de validación. La tesis puede contribuir haciendo explícito el contrato de características, guardando artefactos de ejecución, documentando las suposiciones de la recompensa, comparando contra Random Forest bajo particiones emparejadas, y mostrando cómo el rendimiento interno de CICIDS2017 cambia bajo una validación más estricta y restricciones de escala de entrenamiento.

La afirmación final defendible no es que QRDQN sea universalmente mejor que los modelos de NIDS supervisados, ni que el proyecto implemente defensa cibernética autónoma. La afirmación defendible es que la tesis evalúa un defensor a nivel de flujo basado en QRDQN, sensible al coste, bajo un benchmark reproducible y un pipeline de validación, y utiliza líneas base supervisadas, comprobaciones con control de fugas, análisis de escala de entrenamiento e inferencia sobre tráfico de laboratorio para determinar hasta qué punto puede confiarse en esa formulación. Si experimentos posteriores muestran un rendimiento débil entre archivos o con tráfico de laboratorio, eso no invalida la tesis. Agudiza su contribución al demostrar los límites del RL entrenado sobre benchmark bajo un cambio de distribución más realista.
